% GENERATED FILE. Edit canonical lessons, then run npm run generate:books.
% canonical-source-hash: fnv1a64:df4d9cf16b0f5d0f
% canonical-lessons: KA-C216-yasassu, KA-C216-svatantrya, KA-C216-nyaya, KA-C216-javabdari, KA-C216-kartavya

\chapter{Success, Freedom, Justice, Responsibility, Duty}
\label{ch:ka-success-freedom-justice-responsibility-duty}
\hlchaptermodality{216}

\begin{chapteropening}
\textbf{By the end of this chapter:} \emph{I can say success, freedom, justice, responsibility and a duty.}

Complete the last lesson of chapter 216: I can say success, freedom, justice, responsibility and a duty.
\end{chapteropening}

\section[\texorpdfstring{\kn{ಯಶಸ್ಸು} (yaśassu)}{ಯಶಸ್ಸು (yaśassu)}]{\kn{ಯಶಸ್ಸು} (yaśassu) — success}

\label{lesson:KA-C216-yasassu}



\begin{quote}
Before the new one: say the Kannada for a tradition, then the Kannada for culture.
\end{quote}

\subsection*{You'll want to know: \kn{ಯಶಸ್ಸು}}
\textbf{\kn{ಯಶಸ್ಸು}} — \emph{yaśassu} — \textquotedblleft{}success\textquotedblright{}.

\begin{grammarlens}[title={What you've built}]
One more word for this chapter.
\end{grammarlens}

\subsection*{Your turn}
\begin{itemize}
\raggedright
  \item \emph{Say it:} \emph{yaśassu}
  \item \emph{Say it:} \emph{yaśassu}, once more
  \item \emph{Say it:} say \emph{yaśassu}
  \item \emph{Recall:} say the Kannada for a tradition, then the Kannada for culture, then say \emph{yaśassu} again
\end{itemize}

\begin{tcolorbox}[breakable,colback=teal!4,colframe=teal!35!black,title={Before you move on}]
What does \kn{ಯಶಸ್ಸು} mean? (\textquotedblleft{}Success\textquotedblright{}.) Say it once more.
\end{tcolorbox}
\section[\texorpdfstring{\kn{ಸ್ವಾತಂತ್ರ್ಯ} (svātantrya)}{ಸ್ವಾತಂತ್ರ್ಯ (svātantrya)}]{\kn{ಸ್ವಾತಂತ್ರ್ಯ} (svātantrya) — freedom}

\label{lesson:KA-C216-svatantrya}



\begin{quote}
Before the new one: say the Kannada for culture, then the Kannada for success.
\end{quote}

\subsection*{You'll want to know: \kn{ಸ್ವಾತಂತ್ರ್ಯ}}
\textbf{\kn{ಸ್ವಾತಂತ್ರ್ಯ}} — \emph{svātantrya} — \textquotedblleft{}freedom\textquotedblright{}.

\begin{grammarlens}[title={What you've built}]
One more word for this chapter.
\end{grammarlens}

\subsection*{Your turn}
\begin{itemize}
\raggedright
  \item \emph{Say it:} \emph{svātantrya}
  \item \emph{Say it:} \emph{svātantrya}, once more
  \item \emph{Say it:} say \emph{svātantrya}
  \item \emph{Recall:} say the Kannada for culture, then the Kannada for success, then say \emph{svātantrya} again
\end{itemize}

\begin{tcolorbox}[breakable,colback=teal!4,colframe=teal!35!black,title={Before you move on}]
What does \kn{ಸ್ವಾತಂತ್ರ್ಯ} mean? (\textquotedblleft{}Freedom\textquotedblright{}.) Say it once more.
\end{tcolorbox}
\section[\texorpdfstring{\kn{ನ್ಯಾಯ} (nyāya)}{ನ್ಯಾಯ (nyāya)}]{\kn{ನ್ಯಾಯ} (nyāya) — justice, fairness}

\label{lesson:KA-C216-nyaya}



\begin{quote}
Before the new one: say the Kannada for success, then the Kannada for freedom.
\end{quote}

\subsection*{You'll want to know: \kn{ನ್ಯಾಯ}}
\textbf{\kn{ನ್ಯಾಯ}} — \emph{nyāya} — \textquotedblleft{}justice, fairness\textquotedblright{}.

\begin{grammarlens}[title={What you've built}]
One more word for this chapter.
\end{grammarlens}

\subsection*{Your turn}
\begin{itemize}
\raggedright
  \item \emph{Say it:} \emph{nyāya}
  \item \emph{Say it:} \emph{nyāya}, once more
  \item \emph{Say it:} say \emph{nyāya}
  \item \emph{Recall:} say the Kannada for success, then the Kannada for freedom, then say \emph{nyāya} again
\end{itemize}

\begin{tcolorbox}[breakable,colback=teal!4,colframe=teal!35!black,title={Before you move on}]
What does \kn{ನ್ಯಾಯ} mean? (\textquotedblleft{}Justice, fairness\textquotedblright{}.) Say it once more.
\end{tcolorbox}
\section[\texorpdfstring{\kn{ಜವಾಬ್ದಾರಿ} (javābdāri)}{ಜವಾಬ್ದಾರಿ (javābdāri)}]{\kn{ಜವಾಬ್ದಾರಿ} (javābdāri) — responsibility}

\label{lesson:KA-C216-javabdari}



\begin{quote}
Before the new one: say the Kannada for freedom, then the Kannada for justice.
\end{quote}

\subsection*{You'll want to know: \kn{ಜವಾಬ್ದಾರಿ}}
\textbf{\kn{ಜವಾಬ್ದಾರಿ}} — \emph{javābdāri} — \textquotedblleft{}responsibility\textquotedblright{}.

\begin{grammarlens}[title={What you've built}]
One more word for this chapter.
\end{grammarlens}

\subsection*{Your turn}
\begin{itemize}
\raggedright
  \item \emph{Say it:} \emph{javābdāri}
  \item \emph{Say it:} \emph{javābdāri}, once more
  \item \emph{Say it:} say \emph{javābdāri}
  \item \emph{Recall:} say the Kannada for freedom, then the Kannada for justice, then say \emph{javābdāri} again
\end{itemize}

\begin{tcolorbox}[breakable,colback=teal!4,colframe=teal!35!black,title={Before you move on}]
What does \kn{ಜವಾಬ್ದಾರಿ} mean? (\textquotedblleft{}Responsibility\textquotedblright{}.) Say it once more.
\end{tcolorbox}
\section[\texorpdfstring{\kn{ಕರ್ತವ್ಯ} (kartavya)}{ಕರ್ತವ್ಯ (kartavya)}]{\kn{ಕರ್ತವ್ಯ} (kartavya) — a duty}

\label{lesson:KA-C216-kartavya}



\begin{quote}
Before the new one: say the Kannada for justice, then the Kannada for responsibility.
\end{quote}

\subsection*{You'll want to know: \kn{ಕರ್ತವ್ಯ}}
\textbf{\kn{ಕರ್ತವ್ಯ}} — \emph{kartavya} — \textquotedblleft{}a duty\textquotedblright{}.

\begin{grammarlens}[title={What you've built}]
One more word for this chapter.
\end{grammarlens}

\subsection*{Your turn}
\begin{itemize}
\raggedright
  \item \emph{Say it:} \emph{kartavya}
  \item \emph{Say it:} \emph{kartavya}, once more
  \item \emph{Say it:} say \emph{kartavya}
  \item \emph{Recall:} say the Kannada for justice, then the Kannada for responsibility, then say \emph{kartavya} again
\end{itemize}

\begin{tcolorbox}[breakable,colback=teal!4,colframe=teal!35!black,title={Before you move on}]
What does \kn{ಕರ್ತವ್ಯ} mean? (\textquotedblleft{}A duty\textquotedblright{}.) Say it once more.
\end{tcolorbox}
